% Lösungen Einheit 1 von 3 – Das Intervall lesen und deuten: die Überdeckungsdeutung (verträglich heißt: der angenommene Anteil wird vom Intervall überdeckt (zusammenbau v0.1)
\einheitenkopf{Einheit 1 von 3 · Das Intervall lesen und deuten: die Überdeckungsdeutung (verträglich heißt: der angenommene Anteil wird vom Intervall überdeckt}

\erg{9}{a) richtig gedeutet \quad b) etwa $[0{,}22; 0{,}40]$ – Waagerechte bei $h = 0{,}30$ \quad c) etwa $[0{,}33; 0{,}52]$; es überdeckt $0{,}50$ – die Vermutung ist verträglich \quad d) etwa 51\,\% (jeder Anteil über 50\,\% bis 52\,\%, ebenso ab 43\,\% bis unter 44\,\%); eine Senkrechte dort schneidet genau 6 Intervalle \quad e) Y ist binomialverteilt mit $n = 25$ und Trefferwahrscheinlichkeit $0{,}9$: jedes Intervall überdeckt p unabhängig mit $0{,}9$. $P(Y = 23) = 300 \cdot 0{,}9^{23} \cdot 0{,}1^2 \approx 0{,}266$ – kleiner als $0{,}27$}
\erg{10}{a) Fehler: Senkrechte bei p statt Waagerechte beim Stichprobenanteil. Richtig: etwa $[0{,}67; 0{,}89]$ aus der Waagerechten bei $h = 0{,}8$ \quad b) Jede Stichprobe liefert unabhängig ein Intervall, das den festen Anteil überdeckt oder nicht – ein Bernoulli-Versuch; die Trefferwahrscheinlichkeit ist die Sicherheitswahrscheinlichkeit}
